%% ---------------------------------------------------------------------------
\section{Introduction}
\label{sec:intro}

Content provenance has left the proposal stage. Cameras sign at
capture~\cite{cai2023leica,sony2024c2pa}, phones attach content credentials of
their own~\cite{samsung2025s25,google2025pixel}, and editing tools record what
they changed in a signed C2PA manifest that travels inside the file. Each
credential rests on a COSE signature over a claim about the content and on an
X.509 certificate path from the signer to an anchor on the C2PA trust list,
and every signature in it, on the claim, on the certificates and on the
timestamp, is RSA or ECDSA. A cryptographically relevant quantum computer
breaks all of them~\cite{shor1997}, and the PKI being built for credentials
today will still be in service when one arrives: the roots on the C2PA trust list are issued for
a median of \NtlRootYearsMed{} years, most of them into the 2040s and 2050s.

The replacements exist. NIST standardised ML-DSA and SLH-DSA in
2024~\cite{fips204,fips205}; RFC~9964 assigned the COSE identifiers that let an
ML-DSA signature appear in a manifest~\cite{rfc9964} and RFC~9881 those for
ML-DSA certificates~\cite{rfc9881}; and in October 2026 version~2.5 of the
C2PA specification admitted ML-DSA-44, 65 and 87~\cite{c2pa25}. The reference
implementation had not followed by the time of writing~\cite{c2pars}. To our
knowledge nobody has measured what the change does to a credential, and the published
discussion of post-quantum image authenticity proposes new
schemes~\cite{iavich2025pqveritas,alkhatib2025deepfake} rather than pricing the
one that ships.

\paragraph{Research question}
This paper asks what admitting ML-DSA costs a conformant C2PA content
credential, where that cost sits, which choices the specification already
allows reduce it without weakening what a verifier can establish, and what a
migrated credential still exposes. We break the question into three parts.
\begin{itemize}
\item[\textbf{RQ1}] How large and how slow is a conformant post-quantum
credential, timestamped and with revocation evidence, and what does it add to
the files that signers actually publish?
\item[\textbf{RQ2}] Which carriage, trust-list and verifier choices recover
that cost, on servers and on constrained devices, and what does each give up?
\item[\textbf{RQ3}] What does a partly or fully migrated credential still
expose: which partial migrations leave which properties forgeable, and does the
reference implementation enforce what a migration relies on, that time-stamping keys cannot attribute content, that
timestamps are judged as the specification says, and that revoked signers stay
revoked?
\end{itemize}

Size is a security question here, not only a performance one. A credential
that outgrows the space a generator reserves forces a second timestamp request
on every capture; a verifier that cannot hold the chain is pushed towards
thumbprints, which an offline verifier cannot resolve; and a migration that
moves the leaf but not its issuer leaves the content forgeable, while one that
moves the signer but not its time-stamping authority leaves the signing time
forgeable, each a distinct property (Table~\ref{tab:properties}). Which of the
cheaper options stay sound is what the three questions together decide. We measure
application processors, a single-board computer and microcontrollers; the
secure elements in which phones sign are outside what we could measure, and
Section~\ref{sec:constrained} discusses them.

\paragraph{What is known, and what is new}
That certificate chains, not handshake signatures, dominate post-quantum
authentication is established for TLS~\cite{sikeridis2020ndss,paquin2020,%
sosnowski2023,westerbaan2021}, and the TLS remedies avoid sending or
signing what the verifier could already hold or
check~\cite{kampanakis2022scas,certabridge,merkletreecerts,schwabe2020kemtls}.
We confirm the structural result for manifests and do not claim it. What is
new is what changes when the signed object is a file rather than a connection.
A manifest pays for its chain once per object and keeps it for the life of the
file. Nothing is negotiated, so the carriage form is a policy every verifier
must accept. Some verifiers are offline for good. And C2PA runs its own trust
lists, whose anchors include issuing CAs as well as roots, so a signer under a
listed issuer can omit its intermediate with no new mechanism, under a reading
of the specification that the reference implementation already accepts and
whose confirmation we have requested. These differences decide which remedies
carry over. We measure them, and what the reference implementation accepts.
The closest analysis of C2PA itself, Golaszewski et al.'s study of
version~2.2~\cite{golaszewski2026eprint}, finds timestamp disagreement, weak
revocation policy and inconsistent validators, but considers neither
post-quantum signatures nor what a credential costs;
Section~\ref{sec:related} sets out what each covers.

\paragraph{Contributions}
\begin{itemize}
\item \emph{The cost, on real files (RQ1).} A measurement of ML-DSA, SLH-DSA,
Falcon and composite credentials against the ES256 and PS256 credentials
signers ship, in bytes, signing and validation time, with and without a
timestamp and stapled revocation evidence. A closed-form model reproduces the
measured manifests of \NmodelN{} schemes to within \NmodelErrMax\%, so the
sizes extend to other schemes by substitution, and substituting our measured
ML-DSA-65 components into \NcpqFiles{} production-signed files from
\NcorpProdGens{} signing applications estimates what the migration adds to
real files (Section~\ref{sec:cost}).
\item \emph{What recovers the cost (RQ2).} Leaf-only carriage for signers and,
to our knowledge for the first time, for time-stamping authorities, checked
against both live C2PA trust lists, the production chains and the reference
implementation, with what it gives up; how a cache of validated signers must
be keyed to stay sound; and what the carriage form costs constrained
verifiers in memory and radio energy, and what each live-video method costs a
constrained signer (Sections~\ref{sec:carrying} and~\ref{sec:edge}).
\item \emph{What a migrated credential exposes, and what the validator
enforces (RQ3).} A migration-order analysis showing which property each
partial migration leaves forgeable (Table~\ref{tab:properties}), and a
conformance matrix of \Nvalcases{} cases through the reference implementation,
patched to sign and validate ML-DSA, including a fully post-quantum credential
with a post-quantum timestamp. The matrix finds two implementation defects,
time-stamping and OCSP-signing certificates accepted as claim signers and a
timestamp the specification says to ignore invalidating the whole manifest,
and one gap between the C2PA profile and RFC~6960: revocation stapled in a
CA-signed OCSP response has no effect (Section~\ref{sec:findings}).
\end{itemize}
Section~\ref{sec:discussion} turns the results into recommendations for CAs,
claim generators, validators and the specification. Every script, raw
measurement, firmware image and patch is released, as is every generated
certificate except the withheld vectors of Section~\ref{sec:validation}, and
every number in the text is generated from the released measurements.
